% Pointwise flat inverse; no standalone preamble.
% Requires amsmath/amssymb/amsthm, theorem environment, and \norm.
% Every label introduced by this fragment has the unique prefix pt-.
\section{Pointwise regularity for arbitrary fine forcing}
\label{pt-section}

The full flat inverse also obeys pointwise estimates. These estimates
concern arbitrary bounded fine forcing, without assuming that it is a
coarse lift. Retain fine spacing $\eta=M^{-1}$, odd $M$, unit coarse
cells, the physical fine measure $\eta^d$, and
\begin{equation}\label{pt-operator}
 \mathcal A_a=d^*d+DRD^*+Q^*aQ,\qquad G_0=\mathcal A_a^{-1},
 \qquad \mathcal L=D^*D.
\end{equation}
Here $d$ is the fine curl, $D$ the scalar gradient, $Q=Q_M^{\rm st}$,
and $R$ is the matched homogeneous flat scalar source projection.
The penalty lies in a fixed interval $0<a_-\le a\le a_+<\infty$;
the normalized source value $a=1$ is included. The discussion applies
to the homogeneous infinite lattice and compatible periodic regulators.

For a fixed coarse-cell patch $U$, define the discrete pointwise norm
\begin{equation}\label{pt-holder-norm}
 \norm{u}_{1,\alpha;U}=\norm{u}_{\infty;U}
  +\sum_\mu\norm{D_\mu^+u}_{\infty;U}
  +\sum_\mu\sup_{\substack{x,y\in U\\x\ne y}}
       \frac{|D_\mu^+u(x)-D_\mu^+u(y)|}{|x-y|_1^\alpha}.
\end{equation}
The points are actual fine-grid points and a fixed collar supplies the
differences. Enlarging the patch by finitely many neighboring coarse
cells changes only constants. Backward differences obey the same bounds.

\begin{theorem}\label{pt-theorem}
For each $0<\alpha<1$ there are $C_\alpha,\rho>0$, independent of
odd $M$ and compatible periodic volume, such that the block from
forcing in coarse cell $m$ to a fixed collar of coarse cell $n$ satisfies
\begin{align}
 \norm{G_0(n,m)}_{\ell^\infty\to C^{1,\alpha}}
       &\le C_\alpha e^{-\rho|n-m|_1},
       \label{pt-pointwise-bound}\\
 \norm{(\mathcal L G_0)(n,m)}_{\ell^\infty\to\ell^\infty}
       &\le C e^{-\rho|n-m|_1}.
       \label{pt-laplacian-bound}
\end{align}
The first target norm is \eqref{pt-holder-norm}; $\mathcal L$ in
the second estimate acts componentwise on the vector output.
On tori use periodic coarse-cell distance. Weighted row and column
block sums are uniformly bounded at every smaller positive exponent.
Consequently bounded fine forcing produces uniformly bounded fields,
first differences, locally $\alpha$-H\"older first differences, and
componentwise Laplacians.
\end{theorem}

The constants need not stay bounded as $\alpha\uparrow1$. In fact a
uniform Lipschitz estimate for the first differences, and a uniform
bound for each individual second difference, are false in dimension
$d\ge2$; an example is given below. No source weak-norm identification
is inferred merely from the theorem's terminology.

\subsection{Reduction to a scalar high-mode inverse}

Use the aliases $p_r=z+2\pi r$, $r\in I_M^d$, with
$I_M=\{-(M-1)/2,\ldots,(M-1)/2\}$, and write
\begin{equation}\label{pt-alias-data}
 u_{r,\mu}^{\pm}=M(e^{\pm ip_{r,\mu}/M}-1),\qquad
 P_r=\sum_\mu u_{r,\mu}^+u_{r,\mu}^-,\qquad
 \mathcal F_M(z)_{sr}=e^{i(z+2\pi r)\cdot s/M}.
\end{equation}
The common holomorphic chart from the full inverse theorem contains
the real fundamental cube in the interior of $|\operatorname{Re}z_\mu|<K<2\pi$.
For real $z$, $\mathcal F_M$ is unitary from alias counting $\ell^2$
to physical-cell $L^2$. On the complex chart it and its inverse have
uniformly bounded $L^2$ norms. In particular a bounded fine-cell input
has alias coefficients bounded in $\ell^2$, and individually, by
$C\norm{f}_\infty$.

Recall the canceled inverse $G_0=F+\beta E_a^{-1}\beta^-$.
Separate the scalar orientation-diagonal multiplier
\begin{equation}\label{pt-high-scalar}
 T_r=P_r^{-1}\quad(r\ne0),\qquad T_0=0.
\end{equation}
The difference $F-T$ is a fixed sum of rank-one operators per vector
orientation: its central block, two mixed blocks, and one high/high
term. Its high vectors are $\omega_{r,\mu}^-/P_r$ and the analytic
transpose partners. The remaining scalar coefficients are uniformly
bounded and holomorphic. Thus the only full-rank pointwise problem is
\eqref{pt-high-scalar}. Minus signs on multipliers throughout this
argument denote evaluation at negative full momentum, with the
analytic transpose; they do not conjugate a complex variable.

\subsection{Weighted heat estimates with complex Bloch twist}

Let $p_s(n)$ be the continuous-time nearest-neighbor heat kernel on
$\mathbb Z$ with generator $f(n+1)+f(n-1)-2f(n)$. Its Poisson-process
representation gives
\begin{equation}\label{pt-heat-moment}
 \sum_n e^{tn}p_s(n)=e^{2s(\cosh t-1)},\qquad
 \sum_n p_s(n)=1.
\end{equation}
It is symmetric and decreasing for $n\ge0$. For completeness, set
$q_n=p_s(n)-p_s(n+1)$ on $n\ge0$. Then
\[
 q_0'=q_1-3q_0,\qquad
 q_n'=q_{n-1}+q_{n+1}-2q_n\quad(n\ge1),
\]
with nonnegative initial data. This bounded cooperative system
preserves nonnegativity, as follows by shifting its generator to a
matrix with nonnegative entries and expanding the exponential.
Also Fourier inversion gives
\begin{equation}\label{pt-heat-origin}
 p_s(0)=\frac1{2\pi}\int_{-\pi}^{\pi}
               e^{-4s\sin^2(t/2)}\,dt
       \le C(1+s)^{-1/2}.
\end{equation}
For $0\le r\le r_0$, summation by parts therefore yields the exact
weighted variation identity
\begin{align}
 \sum_n e^{r|n|}|p_s(n+1)-p_s(n)|
  &=(1+e^r)\left[p_s(0)
       +(1-e^{-r})\sum_{n\ge1}e^{rn}p_s(n)\right]\notag\\
  &\le C\left[(1+s)^{-1/2}
                   +r e^{2s(\cosh r-1)}\right].
       \label{pt-weighted-variation}
\end{align}
Take $s=M^2t$ and $r=\delta/M$, where $\delta$ is a fixed small
constant. Multiplying the difference by $M$ gives
\begin{equation}\label{pt-first-heat}
 \sum_n e^{\delta|n|/M}
       |D_\eta p_{M^2t}(n)|
       \le C(t^{-1/2}+\delta)e^{C\delta^2t}.
\end{equation}
For $t<\eta^2$ the trivial $M$ bound is smaller than $t^{-1/2}$.
The weight is submultiplicative. Apply
\eqref{pt-first-heat} twice to
$p_s=p_{s/2}*p_{s/2}$ to obtain the second-difference bound.
Tensor products then give, for any specified composition of $k$ fine
forward/backward differences in $d$ dimensions,
\begin{equation}\label{pt-heat-two-derivatives}
 \sum_{n\in\mathbb Z^d}e^{\delta|n|_1/M}
       |D_\eta^k p_{M^2t}(n)|
       \le Ct^{-k/2},\qquad 0<t\le1,\quad k=0,1,2.
\end{equation}
The displayed $p$ is a probability kernel. If it is instead written
as a density against physical fine volume, the density factor $M^d$
cancels the summation factor $M^{-d}$.

The reconstructed Bloch heat operator has the exact image-sum kernel
\begin{equation}\label{pt-twisted-heat}
 H_z(t;s,s')=\sum_{k\in\mathbb Z^d}
             e^{-iz\cdot k}p_{M^2t}(s-s'+kM).
\end{equation}
This convention uses ordinary summation over the fine-cell input
index $s'$. The formula follows by finite Fourier inversion and
continues to complex $z$ by the exponential moment bound.
For $s$ in a fixed output collar and $s'$ in the input cell,
\[
 |e^{-iz\cdot k}|
       \le C e^{\delta|s-s'+kM|_1/M}.
\]
Thus \eqref{pt-heat-two-derivatives}, including the corresponding
output differences of the image sum, proves
\begin{equation}\label{pt-twisted-derivatives}
 \norm{D_\eta^kH_z(t)}_{\infty\to\infty}
       \le Ct^{-k/2},\qquad 0<t\le1,\quad k=0,1,2.
\end{equation}
The difference between first-gradient rows at fine points $x,y$ is
bounded both by $Ct^{-1/2}$ and, by telescoping, by
$C|x-y|_1t^{-1}$. Interpolating these two elementary bounds gives
\begin{equation}\label{pt-heat-holder}
 \norm{H_z(t)}_{\infty\to C^{1,\alpha}}
       \le C_\alpha t^{-(1+\alpha)/2},\qquad 0<t\le1.
\end{equation}
Subtracting the central heat mode changes this only by a bounded
term, since $P_0,u_0$ and its plane-wave profile are uniformly bounded
on the chart. For $t\ge1$, the high-alias bound
$\operatorname{Re}P_r\ge c|r|^2$ and the individual input-coefficient
bound instead give
\begin{align}
 &\norm{H_z(t)-e^{-tP_0}\Pi_0(z)}_{\infty\to C^{1,\alpha}}
       \notag\\
 &\qquad\le C\sum_{r\ne0}(1+|r|)^{1+\alpha}e^{-ct|r|^2}
       \le C_\alpha e^{-ct/2}.
       \label{pt-heat-large-time}
\end{align}
The high-mode inverse is therefore the convergent operator integral
\begin{equation}\label{pt-scalar-integral}
 T(z)=\int_0^\infty
             \bigl[H_z(t)-e^{-tP_0}\Pi_0(z)\bigr],dt.
\end{equation}
The small-time integral is finite because
$\int_0^1t^{-(1+\alpha)/2}\,dt=2/(1-\alpha)$.
This proves a uniform holomorphic
$\ell^\infty\to C^{1,\alpha}$ bound for $T$ for every $\alpha<1$.

\subsection{Pointwise bounds for the rank corrections}

The centered scalar and straight-vector averaging multipliers satisfy
\begin{equation}\label{pt-alias-weights}
 |a_r^\pm|\le Cw(r),\qquad
 |\omega_{r,\mu}^\pm|\le\frac{Cw(r)}{1+|r_\mu|},\qquad
 w(r)=\prod_\nu(1+|r_\nu|)^{-1}.
\end{equation}
Indeed, for a nonzero component, the numerator of
$\sin((z_\nu+2\pi r_\nu)/2)/
 [M\sin((z_\nu+2\pi r_\nu)/(2M))]$
is uniformly bounded and its denominator has modulus at least
$c|r_\nu|$ on the enlarged chart. The zero component is bounded
directly. The vector multiplier has one extra longitudinal factor.
This includes aliases at the edge of the fine zone.

On a dyadic maximum-norm shell,
\begin{equation}\label{pt-shell-sum}
 \sum_{2^j\le|r|_\infty<2^{j+1}}w(r)
                      \le C_d(j+1)^{d-1}.
\end{equation}
Choose a coordinate at the maximal scale. Its reciprocal weight sums
to $O(1)$ over that shell, and each other coordinate contributes
$O(j+1)$. Summing over the possible maximal coordinates proves
\eqref{pt-shell-sum}. Consequently
\begin{equation}\label{pt-rank-holder-sum}
 \sum_{r\ne0}(1+|r|)^{1+\alpha}
             \left|\frac{\omega_{r,\mu}^-}{P_r}\right|
   \le C\sum_{r\ne0}w(r)(1+|r|)^{\alpha-1}
   \le C_\alpha.
\end{equation}
Absolute Fourier summation now bounds the reconstructed high rank
vectors in $C^{1,\alpha}$. Their central partners have the same bound.
The Woodbury vector has
\begin{equation}\label{pt-beta-bound}
 |\beta_{r,\mu}|
  \le C\left(\frac{|a_r^-|}{(1+|r|)^3}
                  +\frac{|\omega_{r,\mu}^-|}{(1+|r|)^2}\right),
       \qquad r\ne0,
\end{equation}
so its output profile satisfies the same weighted sum. Its canceled
central component is bounded directly.

The input rows of every rank term have uniformly bounded alias
$\ell^2$ norm. By the physical-cell Fourier normalization,
\[
 |\text{input row}(f)|
       \le C\norm{f}_{L^2(\text{physical cell})}
       \le C\norm{f}_\infty.
\]
The cell has physical volume one; there is no factor $M^{d/2}$.
Combining these row estimates with
\eqref{pt-rank-holder-sum}--\eqref{pt-beta-bound} controls all rank
corrections. The full reconstructed $G_0$ therefore has the required
uniform holomorphic pointwise bound on the common chart.

\subsection{The componentwise Laplacian}

For the scalar term, $\mathcal L T=I-\Pi_0$ is bounded on
$\ell^\infty$. Applying $\mathcal L$ to a high rank vector
$\omega_{r,\mu}^-/P_r$ gives $\omega_{r,\mu}^-$. Its full
reconstruction is a column of the physical adjoint average, minus
its bounded central mode. Explicitly, for $x=n+s/M$, $M=2h+1$,
\begin{equation}\label{pt-adjoint-average}
 (Q^*B)_\mu(x)=b_s B_\mu(n)+(1-b_s)B_\mu(n-e_\mu),
 \qquad b_s=\frac{s_\mu+h+1}{M}.
\end{equation}
These weights are nonnegative and sum to one. The Bloch profile is
$b_s+(1-b_s)e^{-iz_\mu}$, uniformly bounded on the complex chart
and any fixed collar. An absolute alias sum of $\omega$ is unnecessary
and would introduce a spurious logarithmic loss.

For the Woodbury vector, its high-alias formula gives
\begin{equation}\label{pt-beta-laplacian}
 (\mathcal L\beta)_{r,\mu}
  =\frac{P_0u_{r,\mu}^+a_r^-}{P_rD_{s,2}}
           -\omega_{r,\mu}^-k_\mu d_\mu^+.
\end{equation}
The first term is $O(w(r)/(1+|r|))$, absolutely summable by
\eqref{pt-shell-sum}. The second is a bounded coefficient times
the profile \eqref{pt-adjoint-average}, with its central subtraction.
The canceled central Laplacian component is bounded. The same input
row estimates as before now prove a uniform holomorphic
$\ell^\infty\to\ell^\infty$ bound for $\mathcal L G_0$.
This argument supplies neither H\"older continuity of its Laplacian
nor bounds on all separate second differences. In particular the
$Q^*$ profile can have transverse block jumps.

\subsection{Gluing, exponential decay, and bounded forcing}

The whole reconstructed inverse and its fine derivatives are periodic
in coarse momentum. The full real symbol is invariant under the
alias relabeling induced by $z\mapsto z+2\pi n$; the within-cell
integer phase cancels wrapping of the alias index. The enlarged
holomorphic charts agree on their real overlaps and hence glue.
Individual separated central/high terms need not be periodic; they
were used only to bound the whole symbol on each chart.

Output collars consist of finitely many neighboring-cell copies,
with bounded holomorphic factors $e^{iz\cdot n}$. Thus the common
strip bounds hold as operator bounds into the discrete norm
\eqref{pt-holder-norm}. For each fixed $M$ these are finite-dimensional
Banach spaces, so entrywise holomorphy and the proved uniform norms
give Banach-valued contour shifting. This proves
\eqref{pt-pointwise-bound}--\eqref{pt-laplacian-bound} without a
growing-dimension factor.

In infinite volume the resulting absolutely summable blocks define
the inverse on bounded forcing and agree with the Hilbert inverse
on finitely supported fine forcing. The inverse identity extends
through absolute convergence of the exponentially summable operator
products at each fixed $M$. The source projection is not local, but
$R=\mathcal L Z\mathcal L$ has exponentially summable cell $L^2$
blocks. Finite-cell norm equivalence makes their pointwise blocks
summable for each fixed $M$. That last identity argument need not
be uniform in $M$: the pointwise inverse bound was proved uniformly
by the heat and rank estimates above.

Finite compatible tori sample the same symbol, giving exact
periodizations and the same bounds in periodic distance, including
short periods and multiplicities. Summing the cell kernels bounds
the field and gradient globally and the gradient H\"older seminorm
on every fixed collar. Nearby points lie in one enlarged collar;
farther points use twice the gradient supremum. This also gives the
global discrete H\"older bound, with periodic physical distance on
tori.

\subsection{A sharp obstruction to the Lipschitz endpoint}

For $d\ge2$, take the compatible torus consisting of one coarse
cell. At its sole coarse momentum, zero, the full inverse is
$a^{-1}$ on the central constant mode and $P_r^{-1}$ on every high
alias. Thus on fine mean-zero forcing it is exactly the scalar
periodic inverse Laplacian, independently in each vector component.

Choose a smooth cutoff $\chi$ supported in a small ball inside the
unit continuum torus and equal to one near its center. Set
\begin{equation}\label{pt-endpoint-example}
 u_\varepsilon(x)=\chi(x)x_1x_2\,
                       \frac12\log(|x|^2+\varepsilon^2).
\end{equation}
For $P=x_1x_2$ and $f(r)=\frac12\log(r^2+\varepsilon^2)$,
harmonicity and homogeneity of $P$ give
\begin{equation}\label{pt-endpoint-laplacian}
 \Delta(Pf)=P\left(f''+\frac{d+3}{r}f'\right)
   =P\frac{(d+2)r^2+(d+4)\varepsilon^2}
                  {(r^2+\varepsilon^2)^2}.
\end{equation}
Since $|P|\le r^2/2$, this is bounded independently of $\varepsilon$.
The cutoff-region derivatives are also uniformly bounded because
that region stays away from zero. However
\begin{equation}\label{pt-endpoint-divergence}
 \partial_1\partial_2u_\varepsilon(0)=\log\varepsilon.
\end{equation}

For fixed $\varepsilon>0$, sample the smooth periodic function on
increasingly fine odd grids and subtract its exact discrete mean.
Define $f_{M,\varepsilon}=\mathcal L_\eta u_{M,\varepsilon}$.
Its mean is zero, and Taylor convergence gives
$\norm{f_{M,\varepsilon}}_\infty\le C+1$ for sufficiently large
$M$, depending on $\varepsilon$. The full inverse returns exactly
$u_{M,\varepsilon}$. Its mixed second fine difference at the origin
converges to \eqref{pt-endpoint-divergence}. Choosing small
$\varepsilon$ and then large $M$ proves that bounded forcing can
produce arbitrarily large mixed second differences. A uniform
Lipschitz bound on first differences would bound these differences
and is therefore impossible. This does not contradict
\eqref{pt-laplacian-bound}, which controls their Laplacian sum rather
than each second difference separately.

\subsection{Source scope}

The theorem supplies arbitrary-forcing pointwise estimates for the
matched homogeneous flat inverse, including its normalized $a=1$
case. It does not by itself identify every numerical factor or term
in the source's weak norms, establish nonflat estimates, handle
general retained-history domains, or control the complete interacting
current carrier. Those applications require exact norm and coordinate
checks. The false endpoint above must not enter such a match through
an implicit assertion that $C^{1,\alpha}$ controls all second derivatives.
